\documentclass[10pt,a4paper]{amsart}
\usepackage{amsmath,amssymb,amsthm,graphicx}
% Retained minimal preamble for the Stacks Project chapter more-algebra.tex (More on Algebra).
% The chapter file pinned-original-source.tex carries only \input{preamble} on its first line, so no
% preamble text lives inside the chapter itself.  The Stacks Project uses one single shared preamble for
% every chapter; the snapshot stages it identically as sources/stacks-chow/preamble.tex and
% sources/stacks-etale-cohomology/preamble.tex.  The block below is copied mechanically, line for line,
% from sources/stacks-chow/preamble.tex lines 14-18 (the Xy-pic setup) and 154-217 (the native theorem
% declarations together with the mathematical macros).  Nothing was added, removed or redefined; the
% wrapper class and the amsmath/amssymb/amsthm/graphicx base packages are supplied by the assembler.
% Chapter: more-algebra.tex -- More on Algebra.
% For commutative diagrams we use Xy-pic
\usepackage[all]{xy}

% We use 2cell for 2-commutative diagrams.
\xyoption{2cell}

% Theorem environments.
%
\theoremstyle{plain}
\newtheorem{theorem}[subsection]{Theorem}
\newtheorem{proposition}[subsection]{Proposition}
\newtheorem{lemma}[subsection]{Lemma}

\theoremstyle{definition}
\newtheorem{definition}[subsection]{Definition}
\newtheorem{example}[subsection]{Example}
\newtheorem{exercise}[subsection]{Exercise}
\newtheorem{situation}[subsection]{Situation}

\theoremstyle{remark}
\newtheorem{remark}[subsection]{Remark}
\newtheorem{remarks}[subsection]{Remarks}

\numberwithin{equation}{subsection}

% Macros
%
\def\lim{\mathop{\mathrm{lim}}\nolimits}
\def\colim{\mathop{\mathrm{colim}}\nolimits}
\def\Spec{\mathop{\mathrm{Spec}}}
\def\Hom{\mathop{\mathrm{Hom}}\nolimits}
\def\Ext{\mathop{\mathrm{Ext}}\nolimits}
\def\SheafHom{\mathop{\mathcal{H}\!\mathit{om}}\nolimits}
\def\SheafExt{\mathop{\mathcal{E}\!\mathit{xt}}\nolimits}
\def\Sch{\mathit{Sch}}
\def\Mor{\mathop{\mathrm{Mor}}\nolimits}
\def\Ob{\mathop{\mathrm{Ob}}\nolimits}
\def\Sh{\mathop{\mathit{Sh}}\nolimits}
\def\NL{\mathop{N\!L}\nolimits}
\def\CH{\mathop{\mathrm{CH}}\nolimits}
\def\proetale{{pro\text{-}\acute{e}tale}}
\def\etale{{\acute{e}tale}}
\def\QCoh{\mathit{QCoh}}
\def\Ker{\mathop{\mathrm{Ker}}}
\def\Im{\mathop{\mathrm{Im}}}
\def\Coker{\mathop{\mathrm{Coker}}}
\def\Coim{\mathop{\mathrm{Coim}}}

% Boxtimes
%
\DeclareMathSymbol{\boxtimes}{\mathbin}{AMSa}{"02}

%
% Macros for moduli stacks/spaces
%
\def\QCohstack{\mathcal{QC}\!\mathit{oh}}
\def\Cohstack{\mathcal{C}\!\mathit{oh}}
\def\Spacesstack{\mathcal{S}\!\mathit{paces}}
\def\Quotfunctor{\mathrm{Quot}}
\def\Hilbfunctor{\mathrm{Hilb}}
\def\Curvesstack{\mathcal{C}\!\mathit{urves}}
\def\Polarizedstack{\mathcal{P}\!\mathit{olarized}}
\def\Complexesstack{\mathcal{C}\!\mathit{omplexes}}
% \Pic is the operator that assigns to X its picard group, usage \Pic(X)
% \Picardstack_{X/B} denotes the Picard stack of X over B
% \Picardfunctor_{X/B} denotes the Picard functor of X over B
\def\Pic{\mathop{\mathrm{Pic}}\nolimits}
\def\Picardstack{\mathcal{P}\!\mathit{ic}}
\def\Picardfunctor{\mathrm{Pic}}
\def\Deformationcategory{\mathcal{D}\!\mathit{ef}}

% --- appended 2026-10-12 (Chief): the project own definitions of the reference/slogan/history
% display environments, copied mechanically from the shared Stacks preamble lines 8-12
% (\usepackage{verbatim} plus the three \newenvironment lines, which the chapter preamble keeps
% verbatim from the project preamble). Nothing added, removed or redefined: these environments
% carry no mathematical content and the project itself typesets them as comments. Without these
% three lines a counted statement whose own text uses them fails with Environment slogan undefined.
\usepackage{verbatim}
\newenvironment{reference}{\comment}{\endcomment}
\newenvironment{slogan}{\comment}{\endcomment}
\newenvironment{history}{\comment}{\endcomment}
\usepackage{cleveref}
\crefname{originallemma}{original lemma}{original lemmas}
\Crefname{originallemma}{Original lemma}{Original lemmas}
\makeatletter
\newlabel{algebra-lemma-disjoint-implies-product}{{00EM}{source}{}{sourcepointer-1}{}}
\newlabel{algebra-lemma-disjoint-implies-product@cref}{{[originallemma][0][]00EM}{[1][0][]source}{}{}{}}
\makeatother
\crefname{originallemma}{original lemma}{original lemmas}
\Crefname{originallemma}{Original lemma}{Original lemmas}
\makeatletter
\newlabel{categories-definition-invertible}{{0FFN}{source}{}{sourcepointer-2}{}}
\newlabel{categories-definition-invertible@cref}{{[originallemma][0][]0FFN}{[1][0][]source}{}{}{}}
\makeatother
\crefname{originallemma}{original lemma}{original lemmas}
\Crefname{originallemma}{Original lemma}{Original lemmas}
\makeatletter
\newlabel{lemma-have-dual-derived}{{38581}{source}{}{sourcepointer-3}{}}
\newlabel{lemma-have-dual-derived@cref}{{[originallemma][0][]38581}{[1][0][]source}{}{}{}}
\makeatother
\crefname{originallemma}{original lemma}{original lemmas}
\Crefname{originallemma}{Original lemma}{Original lemmas}
\makeatletter
\newlabel{lemma-lift-perfect-from-residue-field}{{19594}{source}{}{sourcepointer-4}{}}
\newlabel{lemma-lift-perfect-from-residue-field@cref}{{[originallemma][0][]19594}{[1][0][]source}{}{}{}}
\makeatother
\title{An object of the derived category is invertible exactly when it is locally a shift of the ring, and then its cohomology is finite projective in a single degree per factor}
\author{The Stacks Project Authors}
\date{stacks-project:more-algebra.tex:local-snapshot-20261012; GNU FDL 1.2 or later}
\begin{document}
\maketitle
\setcounter{subsection}{1}
\noindent\textit{Editorial note.} All delivered mathematical text is the author's own, copied from the pinned Stacks Project chapter file more-algebra.tex; no statement or proof step was rewritten, reordered or removed.\par\medskip
\noindent\textit{Editorial note.} The chapter file pinned-original-source.tex carries only \textbackslash{}input\{preamble\} on its first line, so no preamble text lives inside the chapter. The retained file sources/stacks-more-algebra/retained-preamble.tex is a mechanical, line-for-line copy of the project's single shared preamble as staged in the snapshot; nothing was added, removed or redefined and no package was added.\par\medskip
\noindent\textit{Editorial note.} The source URL is the source-level Stacks Project URL, as required when more than one card is produced from one source; the per-claim tag is carried by claim\_key.\par\medskip
\noindent\textit{Editorial note.} The author writes 'we omit the details' at proof line 38682 for the local checking of the composition map; the phrase is kept verbatim. The four references inside the counted proof span are recorded as non-load-bearing navigation pointers.\par\medskip
\noindent\textit{Allowed standard prerequisites (editorial).}\begin{itemize}
\item Invertible objects in a symmetric monoidal category and their left duals. condition (1) at statement lines 38661-38662 and used at proof lines 38683 and 38686-38688. Is the notion at issue and the source of the left dual that gives perfection.
\item The residue-field fibre M tensor\textasciicircum{}L\_R kappa and its cohomology dimensions. used at proof lines 38689-38693. Forces the single-degree concentration that characterises condition (2).
\item Tor amplitude of a perfect complex and the product decomposition of R over disjoint open subsets of the spectrum. used at proof lines 38701-38705. Produces the global product decomposition R = prod R\_n and hence the remaining assertions.
\end{itemize}\medskip
\ifdefined\hypertarget\hypertarget{sourcepointer-1}{}\fi
\noindent\textit{Source reference.} Original lemma 00EM: focused-seeds/AGP/sources/stacks-snapshot-20261011/tags.txt, source lines 540--540. Cross-chapter navigation pointer into the Stacks Project; the standard fact it supplies is applied inside the counted proof span and the argument resting on it stays there.\par\medskip
\ifdefined\hypertarget\hypertarget{sourcepointer-2}{}\fi
\noindent\textit{Source reference.} Original lemma 0FFN: focused-seeds/AGP/sources/stacks-snapshot-20261011/tags.txt, source lines 19293--19293. Cross-chapter navigation pointer into the Stacks Project; the standard fact it supplies is applied inside the counted proof span and the argument resting on it stays there.\par\medskip
\ifdefined\hypertarget\hypertarget{sourcepointer-3}{}\fi
\noindent\textit{Source reference.} Original lemma 38581: sources/stacks-more-algebra/pinned-original-source.tex, source lines 38581--38581. Same-chapter navigation pointer; the statement it names is standard for the section and the argument that consumes it stays inside the counted proof span.\par\medskip
\ifdefined\hypertarget\hypertarget{sourcepointer-4}{}\fi
\noindent\textit{Source reference.} Original lemma 19594: sources/stacks-more-algebra/pinned-original-source.tex, source lines 19594--19594. Same-chapter navigation pointer; the statement it names is standard for the section and the argument that consumes it stays inside the counted proof span.\par\medskip
\begin{lemma}
\label{lemma-invertible-derived}
Let $R$ be a ring. Let $M$ be an object of $D(R)$. The following
are equivalent
\begin{enumerate}
\item $M$ is invertible in $D(R)$, see
Categories, Definition \ref{categories-definition-invertible}, and
\item for every prime ideal $\mathfrak p \subset R$ there
exists an $f \in R$, $f \not \in \mathfrak p$ such that
$M_f \cong R_f[-n]$ for some $n \in \mathbf{Z}$.
\end{enumerate}
Moreover, in this case
\begin{enumerate}
\item[(a)] $M$ is a perfect object of $D(R)$,
\item[(b)] $M = \bigoplus H^n(M)[-n]$ in $D(R)$,
\item[(c)] each $H^n(M)$ is a finite projective $R$-module,
\item[(d)] we can write $R = \prod_{a \leq n \leq b} R_n$
such that $H^n(M)$ corresponds to an invertible $R_n$-module.
\end{enumerate}
\end{lemma}

\begin{proof}
Assume (2). Consider the object $R\Hom_R(M, R)$ and the composition map
$$
R\Hom(M, R) \otimes_R^\mathbf{L} M \to R
$$
Checking locally we see that this is an isomorphism; we omit the details.
Because $D(R)$ is symmetric monoidal we see that $M$ is invertible.

\medskip\noindent
Assume (1). Observe that an invertible object of a monoidal category
has a left dual, namely, its inverse. Thus $M$ is perfect by
Lemma \ref{lemma-have-dual-derived}. Consider a prime ideal
$\mathfrak p \subset R$ with residue field $\kappa$.
Then we see that $M \otimes_R^\mathbf{L} \kappa$ is an
invertible object of $D(\kappa)$. Clearly this implies that
$\dim H^i(M \otimes_R^\mathbf{L} \kappa)$ is nonzero exactly
for one $i$ and equal to $1$ in that case. By
Lemma \ref{lemma-lift-perfect-from-residue-field}
this gives (2).

\medskip\noindent
In the proof above we have seen that (a) holds. Let $U_n \subset \Spec(R)$
be the union of the opens of the form $D(f)$ such that
$M_f \cong R_f[-n]$. Clearly, $U_n \cap U_{n'} = \emptyset$ if $n \not = n'$.
If $M$ has tor amplitude in $[a, b]$, then $U_n = \emptyset$
if $n \not \in [a, b]$. Hence we see that we have a product decomposition
$R = \prod_{a \leq n \leq b} R_n$ as in (d) such that
$U_n$ corresponds to $\Spec(R_n)$, see
Algebra, Lemma \ref{algebra-lemma-disjoint-implies-product}.
Since $D(R) = \prod_{a \leq n \leq b} D(R_n)$ and similarly
for the category of modules parts (b), (c), and (d) follow immediately.
\end{proof}

\end{document}
